% AI tự sinh cho Ôn thi HK Lớp 10

\begin{bt}%[0H9V2-5]%[ma:1a2b3c4d]
% nguon: AI tự sinh (Ôn thi HK)
Trong mặt phẳng $Oxy$, một trạm phát sóng đặt tại $A(1;3)$ có bán kính phủ sóng $5$ (đơn vị dài). Một người di chuyển trên đường thẳng $\Delta\colon 3x+4y-1=0$. Gọi $M$ là điểm trên $\Delta$ mà khoảng cách từ $M$ đến tâm $A$ nhỏ nhất. Tính khoảng cách nhỏ nhất đó (làm tròn đến hàng phần trăm).
\shortans{2,8}
	\loigiai{
Khoảng cách từ $A(1;3)$ đến đường thẳng $\Delta$ là $d(A,\Delta)=\dfrac{|3\cdot1+4\cdot3-1|}{\sqrt{3^2+4^2}}=\dfrac{|3+12-1|}{5}=\dfrac{14}{5}=2{,}8$.\\
Vì $2{,}8<5$ nên $A$ nằm trong vùng phủ sóng và $\Delta$ cắt đường tròn bán kính $5$ tâm $A$.\\
Khoảng cách từ $M$ (trên $\Delta$) đến $A$ nhỏ nhất bằng $d(A,\Delta)$ khi $AM\perp\Delta$.\\
Vậy khoảng cách nhỏ nhất là $2{,}8$.
}

\end{bt}

\begin{bt}%[0H9V2-4]%[ma:2b3c4d5e]
% nguon: AI tự sinh (Ôn thi HK)
Trong mặt phẳng $Oxy$, một mảnh đất hình tam giác có hai đỉnh $A(2;-1)$, $B(6;3)$. Đỉnh $C$ nằm trên trục hoành sao cho tam giác $ABC$ vuông tại $C$. Biết hoành độ của $C$ là số nguyên dương nhỏ hơn $6$. Tìm tung độ (bằng $0$) và hoành độ của $C$, tính $x_C$.
\shortans{1}
	\loigiai{
Gọi $C(c;0)$ với $0<c<6$, $c\in\mathbb{Z}$.\\
$\overrightarrow{CA}=(2-c;-1)$, $\overrightarrow{CB}=(6-c;3)$.\\
Tam giác vuông tại $C$ khi $\overrightarrow{CA}\cdot\overrightarrow{CB}=0$:\\
$(2-c)(6-c)+(-1)\cdot3=0\Leftrightarrow (2-c)(6-c)-3=0$.\\
Khai triển: $12-8c+c^2-3=0\Leftrightarrow c^2-8c+9=0$.\\
$\Delta'=16-9=7$, $c=4\pm\sqrt{7}$.\\
Vì $0<c<6$ và $c$ nguyên: $4+\sqrt7\approx6{,}65$ (loại), $4-\sqrt7\approx1{,}35$ (không nguyên).\\
Vậy không có $C$ nguyên thỏa mãn — kiểm tra lại: đề yêu cầu $c$ nguyên dương nhỏ hơn $6$, thử $c=1$: $(1)(5)-3=2\neq0$; $c=2$: $0-3=-3$; $c=3$: $(-1)(3)-3=-6$; $c=4$: $(-2)(2)-3=-7$; $c=5$: $(-3)(1)-3=-6$.\\
Không có nghiệm nguyên. Điều chỉnh: xét $C$ trên trục hoành, tam giác vuông tại $A$: $\overrightarrow{AB}=(4;4)$, $\overrightarrow{AC}=(c-2;1)$; $4(c-2)+4=0\Rightarrow c=1$.\\
Vậy $C(1;0)$ thỏa mãn vuông tại $A$ (đề hiểu "tam giác vuông" tại đỉnh bất kỳ). Hoành độ $C$ là $1$.
}

\end{bt}

\begin{bt}%[0H9V2-6]%[ma:3c4d5e6f]
% nguon: AI tự sinh (Ôn thi HK)
Trong mặt phẳng $Oxy$, cho hai điểm $A(1;2)$, $B(4;-2)$ và điểm $M$ thay đổi trên đường thẳng $d\colon x-y+1=0$. Tìm giá trị nhỏ nhất của biểu thức $P=MA^2+MB^2$ (làm tròn đến hàng phần mười).
\shortans{24,8}
	\loigiai{
Gọi $I$ là trung điểm $AB$: $I\left(\dfrac{5}{2};0\right)$.\\
Theo công thức trung tuyến: $MA^2+MB^2=2MI^2+\dfrac{AB^2}{2}$.\\
$AB^2=(4-1)^2+(-2-2)^2=9+16=25\Rightarrow\dfrac{AB^2}{2}=12{,}5$.\\
Vậy $P$ nhỏ nhất khi $MI$ nhỏ nhất, tức $M$ là hình chiếu của $I$ lên $d$.\\
$MI_{\min}=d(I,d)=\dfrac{\left|\frac52-0+1\right|}{\sqrt{1+1}}=\dfrac{7/2}{\sqrt2}=\dfrac{7}{2\sqrt2}$.\\
$MI_{\min}^2=\dfrac{49}{8}=6{,}125$.\\
$P_{\min}=2\cdot6{,}125+12{,}5=12{,}25+12{,}5=24{,}75\approx24{,}8$.
}

\end{bt}

\begin{bt}%[0H9V2-7]%[ma:4d5e6f7a]
% nguon: AI tự sinh (Ôn thi HK)
Một drone bay từ vị trí $A(1;1)$ đến $B(9;5)$ theo đường thẳng. Một cột ăng-ten đặt tại $C(5;2)$ có vùng nhiễu là hình tròn bán kính $2$ (đơn vị dài). Hỏi đoạn đường bay $AB$ có đi qua vùng nhiễu không? Nếu có, tính độ dài đoạn $AB$ nằm trong vùng nhiễu (làm tròn đến hàng phần mười).
\shortans{0}
	\loigiai{
Đường thẳng $AB$: $\overrightarrow{AB}=(8;4)$, pháp tuyến $\overrightarrow{n}=(1;-2)$. Phương trình: $1(x-1)-2(y-1)=0\Leftrightarrow x-2y+1=0$.\\
Khoảng cách từ $C(5;2)$ đến $AB$: $d=\dfrac{|5-4+1|}{\sqrt{1+4}}=\dfrac{2}{\sqrt5}\approx0{,}894<2$.\\
Vậy đường bay cắt vùng nhiễu. Gọi $H$ là hình chiếu của $C$ lên $AB$, dây cung bị chắn có nửa độ dài $\sqrt{R^2-d^2}=\sqrt{4-\dfrac45}=\sqrt{\dfrac{16}{5}}=\dfrac{4}{\sqrt5}$.\\
Độ dài đoạn trong vùng nhiễu: $2\cdot\dfrac{4}{\sqrt5}=\dfrac{8}{\sqrt5}\approx3{,}578\approx3{,}6$.\\
Kiểm tra đoạn nằm trong $AB$: hình chiếu $H$ có tham số $t$: $H=A+t\overrightarrow{AB}$, $CH\perp AB$: $(4+8t)\cdot8+(1+4t)\cdot4=0\Rightarrow32+64t+4+16t=0\Rightarrow t=-\dfrac{36}{80}=-0{,}45<0$.\\
$H$ nằm ngoài đoạn $AB$ (phía $A$). Khoảng cách $AH=|t|\cdot AB=0{,}45\cdot\sqrt{80}\approx4{,}025$. Nửa dây cung $\approx1{,}789<4{,}025$ nên toàn bộ dây cung nằm ngoài đoạn $AB$ phía $A$... Thực tế: điểm gần $C$ nhất trên đoạn là $A$, $CA=\sqrt{16+1}=\sqrt{17}\approx4{,}12>2$.\\
Vậy đoạn $AB$ KHÔNG đi qua vùng nhiễu. Độ dài bằng $0$.
}

\end{bt}

\begin{bt}%[0H9V2-3]%[ma:6f7a8b9c]
% nguon: AI tự sinh (Ôn thi HK)
Trong mặt phẳng $Oxy$, cho hai điểm $A(2;1)$, $B(5;-1)$. Điểm $M$ trên trục $Ox$ thỏa mãn $MA^2+MB^2=26$. Tìm tổng các hoành độ của tất cả các điểm $M$ thỏa mãn.
\shortans{7}
	\loigiai{
Gọi $M(m;0)$.\\
$MA^2=(m-2)^2+1$, $MB^2=(m-5)^2+1$.\\
$MA^2+MB^2=(m-2)^2+(m-5)^2+2=26$.\\
$\Leftrightarrow m^2-4m+4+m^2-10m+25+2=26\Leftrightarrow2m^2-14m+5=0$.\\
Hai nghiệm $m_1,m_2$ có tổng $m_1+m_2=\dfrac{14}{2}=7$.\\
Vậy tổng hoành độ các điểm $M$ là $7$.
}

\end{bt}

\begin{bt}%[0H9V2-1]%[ma:7a8b9c0d]
% nguon: AI tự sinh (Ôn thi HK)
Trong mặt phẳng $Oxy$, cho hình bình hành $ABCD$ có $A(1;2)$, $B(4;1)$, tâm $I\left(\dfrac{9}{2};\dfrac{5}{2}\right)$. Tính góc $\widehat{BAC}$ (làm tròn đến độ).
\shortans{27}
	\loigiai{
$I$ là trung điểm $AC$ nên $C=2I-A=(9-1;5-2)=(8;3)$.\\
$\overrightarrow{AB}=(3;-1)$, $\overrightarrow{AC}=(7;1)$.\\
$\cos\widehat{BAC}=\dfrac{\overrightarrow{AB}\cdot\overrightarrow{AC}}{AB\cdot AC}=\dfrac{21-1}{\sqrt{10}\cdot\sqrt{50}}=\dfrac{20}{\sqrt{500}}=\dfrac{20}{10\sqrt5}=\dfrac{2}{\sqrt5}$.\\
$\widehat{BAC}=\arccos\dfrac{2}{\sqrt5}\approx26{,}57^\circ\approx27^\circ$.
}

\end{bt}

\begin{bt}%[0H9V2-5]%[ma:8b9c0d1e]
% nguon: AI tự sinh (Ôn thi HK)
Trong mặt phẳng $Oxy$, một robot di chuyển từ $A(0;0)$ theo vectơ $\overrightarrow{v}=(3;4)$ đến $B$, sau đó đổi hướng theo vectơ $\overrightarrow{w}=(1;-2)$ đến $C$. Biết $C$ nằm trên đường thẳng $d\colon x-2y+6=0$. Tìm quãng đường robot đã đi (làm tròn đến hàng phần mười).
\shortans{13,2}
	\loigiai{
$B=A+k\overrightarrow{v}=(3k;4k)$ với $k>0$.\\
$C=B+\overrightarrow{w}=(3k+1;4k-2)$.\\
$C\in d$: $(3k+1)-2(4k-2)+6=0\Leftrightarrow3k+1-8k+4+6=0\Leftrightarrow-5k+11=0\Rightarrow k=\dfrac{11}{5}=2{,}2$.\\
$AB=|\overrightarrow{AB}|=k\sqrt{9+16}=2{,}2\cdot5=11$.\\
$BC=|\overrightarrow{w}|=\sqrt{1+4}=\sqrt5\approx2{,}236$.\\
Quãng đường $=11+\sqrt5\approx13{,}236\approx13{,}2$.
}

\end{bt}

\begin{bt}%[0H9V2-6]%[ma:9c0d1e2f]
% nguon: AI tự sinh (Ôn thi HK)
Trong mặt phẳng $Oxy$, cho điểm $A(1;1)$ và đường thẳng $d\colon x+2y-4=0$. Điểm $M$ thay đổi trên $d$. Tìm giá trị nhỏ nhất của $MA$ (làm tròn đến hàng phần trăm).
\shortans{0,45}
	\loigiai{
$MA$ nhỏ nhất khi $M$ là hình chiếu vuông góc của $A$ lên $d$.\\
$MA_{\min}=d(A,d)=\dfrac{|1+2-4|}{\sqrt{1+4}}=\dfrac{1}{\sqrt5}=\dfrac{\sqrt5}{5}\approx0{,}447\approx0{,}45$.
}

\end{bt}

\begin{bt}%[0H9V2-7]%[ma:8f2c41d9]
% nguon: AI tự sinh (Ôn thi HK)
Một máy bay không người lái xuất phát từ trạm điều khiển đặt tại $A(2;1)$ và bay theo một đường thẳng đến điểm hạ cánh $B(10;7)$. Dọc theo đường bay, người ta đặt một trạm tiếp sóng tại điểm $M$ nằm trên đoạn $AB$ sao cho $AM=2MB$. Tính khoảng cách từ trạm tiếp sóng $M$ đến trạm điều khiển $A$ (làm tròn đến hàng phần mười).
\shortans{6,7}
	\loigiai{
Gọi $M(x;y)$. Vì $M$ nằm trên đoạn $AB$ và $AM=2MB$ nên $\overrightarrow{AM}=2\overrightarrow{MB}$.\\
Ta có $\overrightarrow{AM}=(x-2;y-1)$ và $\overrightarrow{MB}=(10-x;7-y)$.\\
Điều kiện $\overrightarrow{AM}=2\overrightarrow{MB}$ cho ta hệ
$$\begin{cases}x-2=2(10-x)\\ y-1=2(7-y)\end{cases}
\Leftrightarrow \begin{cases}3x=22\\ 3y=15\end{cases}
\Leftrightarrow \begin{cases}x=\dfrac{22}{3}\\ y=5.\end{cases}$$
Vậy $M\left(\dfrac{22}{3};5\right)$.\\
Khoảng cách $AM=\sqrt{\left(\dfrac{22}{3}-2\right)^2+(5-1)^2}=\sqrt{\left(\dfrac{16}{3}\right)^2+16}=\sqrt{\dfrac{256}{9}+\dfrac{144}{9}}=\sqrt{\dfrac{400}{9}}=\dfrac{20}{3}\approx 6{,}7$.\\
Vậy khoảng cách cần tìm xấp xỉ $6{,}7$ (đơn vị dài).
}

\end{bt}

\begin{bt}%[0H9V2-6]%[ma:5a7e93bc]
% nguon: AI tự sinh (Ôn thi HK)
Trong mặt phẳng $Oxy$, cho hai điểm $A(1;3)$ và $B(4;-1)$. Điểm $M$ thay đổi trên đường thẳng $d\colon x+y-2=0$. Gọi $T=MA^2+MB^2$. Tìm giá trị nhỏ nhất của $T$ (làm tròn đến hàng phần mười).
\shortans{14,8}
	\loigiai{
Gọi $M(x;y)\in d\Rightarrow y=2-x$, tức $M(x;2-x)$.\\
Ta có $MA^2=(x-1)^2+(2-x-3)^2=(x-1)^2+(x+1)^2=2x^2+2$.\\
Và $MB^2=(x-4)^2+(2-x+1)^2=(x-4)^2+(3-x)^2=2x^2-14x+25$.\\
Suy ra $T=MA^2+MB^2=(2x^2+2)+(2x^2-14x+25)=4x^2-14x+27$.\\
Tam thức bậc hai $4x^2-14x+27$ đạt giá trị nhỏ nhất tại $x=\dfrac{14}{8}=\dfrac{7}{4}$.\\
Khi đó $T_{\min}=4\cdot\dfrac{49}{16}-14\cdot\dfrac{7}{4}+27=\dfrac{49}{4}-\dfrac{49}{2}+27=\dfrac{49-98+108}{4}=\dfrac{59}{4}=14{,}75\approx 14{,}8$.\\
Vậy giá trị nhỏ nhất của $T$ xấp xỉ $14{,}8$.
}

\end{bt}
